\section{损失函数决定网络在学习什么}
\label{sec:14-Deep-Learning/02-Losses-and-Backpropagation/01}

网络架构只规定``能够表示哪些函数''——那是一张巨大的候选菜单。训练需要有人告诉优化器，菜单上哪些菜比另一些更好。这个角色归损失函数。

想象你要预测北京三环内的房价，手上有面积、楼层、朝向、房龄四个特征。你搭了一个三层全连接网络，结构合理，激活函数也选对了。训练开始之后，网络输出的数字到底应该靠近什么？靠近成交价本身？靠近成交价的对数？靠近某个分位数？不同的选择会让网络学到完全不同的规律，而且结构完全相同的网络自身无法判断哪个才是你真正想要的。

以下对样本 \((x_i,y_i)\) 记经验风险为

\[
\widehat R(\theta)
=\frac1n\sum_{i=1}^n
\ell(f_\theta(x_i),y_i).
\]

不同的 \(\ell\) 不只是曲线形状不同——它们各自对应不同的条件统计量或概率模型，选错 \(\ell\) 等于问错问题。

\subsection{平方损失追的是条件均值，不是天然正确}

对平方损失 \(\ell(\hat y,y)=\frac12(\hat y-y)^2\)，固定 \(X=x\)，条件风险可以按条件均值拆开：

\[
\begin{aligned}
\E[(Y-a)^2\mid X=x]
&=\E\bigl[(Y-\E[Y\mid X=x]+\E[Y\mid X=x]-a)^2\mid X=x\bigr]\\[2mm]
&=\Var(Y\mid X=x)
+\bigl(a-\E[Y\mid X=x]\bigr)^2.
\end{aligned}
\]

第一项是数据自身在 \(x\) 处的波动，与预测 \(a\) 无关。第二项非负，且只在 \(a=\E[Y\mid X=x]\) 时为零。所以平方损失的最优预测恰好是条件均值。

这看起来很合理——均值不是最自然的汇总方式吗？但合理的前提是条件均值确实能代表你想知道的量。回到房价：一个小区最近成交了十套，九套在 400 万到 500 万之间，但有一套顶层复式卖了 2000 万。条件均值被这一套拉到接近 600 万，而你如果用这个模型给下一个普通三居室估价，你会系统性地高估。

这时候你需要的是条件中位数——一半成交价在它之上、一半在它之下，不被极端值拖着走。绝对损失 \(\ell(\hat y,y)=|\hat y-y|\) 恰好对应这个需求。固定 \(X=x\)，考察

\[
\E[|Y-a|\mid X=x]
\]

对 \(a\) 的次导数：左导数是 \(-P(Y<a\mid X=x)\) 减去 \(+P(Y>a\mid X=x)\) 的某种符号组合，整理后零点满足

\[
P(Y<a\mid X=x)=P(Y>a\mid X=x),
\]

恰好是中位数的定义。所以 MAE 的最优预测是条件中位数。

同一个网络、同一批数据，换一个损失函数就换了一个最优预测。MSE 追均值，MAE 追中位数，分位数损失追指定的分位点。网络没有偏好——损失函数替它做了选择。Huber loss 在小残差处二次、大残差处线性，是对效率和抗离群性的人工折中；但它的阈值本身就是在定义``多大的残差算离群''，这也是一个隐藏的选择。

如果业务真正需要的是完整条件分布——比如既要估计期望损失、又要评估尾部风险——单一 MSE 或 MAE 网络输出一个标量，无法传递分布形状的信息。这时需要分位数损失同时训练多个分位点，或者直接让网络输出分布的参数。

\subsection{分类交叉熵背后是一个完整的概率模型}

二分类情形，记真实阳性条件概率为 \(p=P(Y=1\mid x)\)，模型预测为 \(q\)。binary cross-entropy 这个损失项在单个样本上的条件风险是

\[
-p\log q-(1-p)\log(1-q).
\]

把它拆成两部分看：

\[
\underbrace{-\bigl[p\log p+(1-p)\log(1-p)\bigr]}_{H(p)\text{：真实分布的 Bernoulli 熵}}
+\underbrace{\Bigl[p\log\frac{p}{q}+(1-p)\log\frac{1-p}{1-q}\Bigr]}_{\KL(\operatorname{Bern}(p)\|\operatorname{Bern}(q))}.
\]

第一项 \(H(p)\) 不依赖 \(q\)。第二项是 Bernoulli 分布之间的 KL 散度，非负且只在 \(q=p\) 时为零。所以 cross-entropy 在 \(q=p\) 处唯一最小——它迫使模型学会输出真实的条件概率。

说 cross-entropy 是``分类默认损失''，等于什么都没解释。它是一套完整的 scoring rule：诚实报告你估计的真实概率，在期望上获得最低损失；任何偏离都会受罚。这也是它被称作严格 proper scoring rule 的原因。

互斥多分类时，网络输出 logits \(z\in\R^K\)，经 softmax 转化为概率向量 \(p=\softmax(z)\)，categorical cross-entropy 取 \(-\log p_{y}\)（\(y\) 是真实类别）。实现中不应先手算 softmax 概率再取对数——当某个 logit 极大时，指数运算可能溢出，后续的 log 又会吞掉结构，最终贡献一个无意义的数值。把 logits 直接交给框架内置的 fused loss（如 PyTorch 的 CrossEntropyLoss），它在内部使用 log-sum-exp 稳定公式，见下一单元``训练动力学与数值稳定''一节。

这个联合损失对 logits 的梯度极为干净：对 logits \(z\) 和 one-hot 标签 \(y\)，

\[
\frac{\partial\ell}{\partial z}=p-y.
\]

预测概率减去真实标签——一个简洁到可以一眼检查的信号。这个信号会成为下一节反向传播的起点，穿过多层参数，最终驱动所有权重更新。

\subsection{训练目标、评估指标和行动成本不是同一个东西}

到这里你已经看到了两个层面的分离：网络架构决定``能表示什么''，损失函数决定``优化器在找什么''。还有第三个层面。

准确率是人最关心的数字——但准确率含有不可微的 argmax，不能直接作为训练目标。于是你把 cross-entropy 当作可优化的 surrogate，在验证集上监控准确率。这是合理的，但它埋了一个坑：最低 cross-entropy 的模型未必在某个固定阈值下有最高准确率。cross-entropy 惩罚的是概率失准，准确率只关心 0.5 阈值那一刀切在哪里。一个模型可能概率估计很准但阈值附近的排序稍有波动，另一个模型概率全部挤在 0.49 和 0.51 之间但排序完美——谁更优，取决于你接下来拿概率做什么。

业务层面还可能引入完全不同的成本结构。漏诊和误诊的代价不同，推荐系统的 top-K 召回和 MRR 各有侧重，广告出价依赖精确的点击率估计而非分类准确率。这些下游行动的成本结构不进入训练目标——它们应当在验证数据上结合具体成本选择阈值、校准概率、或重新评估排序质量。

类别不平衡时，直接加 class weights 或使用 weighted cross-entropy，会让少数类错误在训练中更昂贵。但这同时改变了训练目标对应的隐式分布：最优解不再是自然人群中的条件概率，而是某个重新加权后的条件概率。模型输出需要在对真实 prior 重新校准后才能解释为概率。过采样、focal loss 和代价敏感学习也在做同一件事——改变训练目标。不能只看``少数类召回上升了''就满意，还需要追问概率含义是否还成立。

总损失常还包含正则项：

\[
L_{\mathrm{total}}
=L_{\mathrm{data}}+\lambda\Omega(\theta).
\]

训练过程中应分开记录 data loss 和 penalty：否则总损失下降可能全部来自权重收缩，而预测拟合反而变差。batch reduction（sum vs. mean）也会改变 data loss 与 weight decay 的相对尺度，换 batch size 时容易被忽略。

\subsection{用极端情形检查语义与实现}

一个可信的损失函数，在动手训练之前应该先通过几个简单测试：

\begin{itemize}
\tightlist
\item
  完全正确且极其自信时，损失是否接近零或理论下限——例如 \(p=1,q=0.999\)，cross-entropy 应该是一个很小的正数；
\item
  完全错误且极其自信时，损失是否足够大且有限——\(p=0,q=0.999\) 时 cross-entropy 应给出约 6.9（以 e 为底），而不是无穷；
\item
  batch 使用 sum 还是 mean，mask 和 sample weights 的分母分别是什么，换 batch size 后 loss 的量级是否按预期缩放；
\item
  标签的取值范围、输出的 shape 和激活函数的概率假设是否互相匹配——比如二分类最后一层是 sigmoid 还是 softmax(dim=2)，相应的标签 shape 和损失函数是不是配套；
\item
  解析梯度能否与自动微分或有限差分在一个小网络上吻合。
\end{itemize}

标签噪声、删失数据、ranking、结构化输出和不对称决策都可能需要专门的损失设计。不要因为框架内置了某个 loss 就把任务压进它的假设里。

\subsubsection{从统计目标推回损失}

用同一组偏斜回归数据分别训练 MSE、absolute 和 0.9 分位数损失模型，比较三个模型输出的条件均值、中位数和高分位数——这三者在该数据上应当明显不同。再为二分类数据加入正类权重，观察排序质量、概率校准和 0.5 阈值准确率在加权前后如何变化。最后用一句话分别说明：这个任务中训练 surrogate 是什么、验证 metric 是什么、部署阶段的 cost 是什么——三者是三个不同的对象。
